\question{مفاهیم پایه‌ای مجموعه‌ها و اعداد}

\subsection{مجموعه‌های عددی معروف}
مجموعه‌ی اعداد طبیعی را با $\IN$ نشان می‌دهیم:
\[
	\IN = \set{1, 2, 3, \dots}
\]
و اعداد صحیح، گویا، حقیقی و مختلط به ترتیب $\IZ$، $\IQ$، $\IR$ و $\IC$ هستند.
برای یک مجموعه‌ی دلخواه $A$، اندازه‌ی آن $\card{A}$ و متمم آن $\setcomp{A}$ است.

\subsection{دنباله‌ها و کران‌ها}
یک دنباله مانند $\seq{a_n}$ را در نظر بگیرید. سقف و کف آن به‌ترتیب
$\ceil{x}$ و $\floor{x}$ تعریف می‌شوند. همچنین اگر $A \subseteq B$، آن‌گاه داریم
$\card{A} \leq \card{B}$ (علامت نابرابری به شکل $\leqslant$ درآمده است).

\begin{proof}
	فرض کنید $A$ و $B$ دو مجموعه باشند. اگر $x \in A$ آن‌گاه $x \in B$
	(تعریف زیرمجموعه). بنابراین $\card{A} \le \card{B}$.
\end{proof}

\subsection{اجزای یک مسئله}
برای حل بسیاری از مسائل، مراحل زیر را طی می‌کنیم:
\begin{ابجد}
\item خواندن صورت سوال و تشخیص ورودی‌ها
\item یافتن ساختار ریاضی مناسب
\item طراحی الگوریتم
\item پیاده‌سازی و آزمون
\end{ابجد}

گاهی هم از شمارش \lr{Abjad} فارسی استفاده می‌کنیم که در محیط \textbf{ابجد} حروف ابجد به کار می‌روند.